\section{The fixed coordinate model}
Write $\Omega=\F_{23}\cup\{\infty\}$, serializing $\infty$ as coordinate $23$. If $Q$ is the set of nonzero squares modulo $23$, let $\mathcal C\subset\F_2^{24}$ be the span of the 23 supports $(a+Q)\cup\{\infty\}$. A word is represented by the binary integer whose $i$th bit is its $i$th coordinate. The twelve words
\[
\begin{gathered}
\mathtt{85335e,
5606be,
2ca9ae,
11fe26,
0f55e2,
046f3e,}\\
\mathtt{02379f,
0149a6,
00a4d3,
0042de,
003e4a,
001f25}
\end{gathered}
\]
form a basis; the notation is hexadecimal. Row reduction of the defining supports gives this basis. Enumerating its $2^{12}$ sums gives weights $0,8,12,16,24$ with multiplicities $1,759,2576,759,1$. This is the projective Golay code \cite{chapman}. We use the Leech model \cite[Theorem 1.4]{brouwer}
\begin{equation}\label{eq:model}
\La=\left\{x/\sqrt8:\begin{array}{l}x\in\Z^{24},\quad x_i\equiv m\pmod2\text{ for one }m\in\{0,1\},\\
\sum_i x_i\equiv4m\pmod8,\quad ((x_i-m)/2\bmod2)_i\in\mathcal C\end{array}\right\}.
\end{equation}
Thus $\La$ is even, unimodular and has minimum squared norm four. These properties and the lattice itself are inherited, not claimed as new.

The three disjoint octads are
\begin{align*}
B_0&=\{0,1,3,4,7,8,16,23\},\\
B_1&=\{10,12,15,18,19,20,21,22\},\\
B_2&=\{2,5,6,9,11,13,14,17\}.
\end{align*}
Put $w_i=2\one_{B_i}/\sqrt8$ and $P_i(x/\sqrt8)=\sum_{j\in B_i}x_j/4$. The matrices
\[
C=\begin{pmatrix}12&5\\1&12\end{pmatrix},\qquad
J=\begin{pmatrix}1&7\\3&22\end{pmatrix},\qquad
T=\begin{pmatrix}4&11\\16&4\end{pmatrix}
\]
act on $\Omega$ by fractional linear transformations over $\F_{23}$ and permute lattice coordinates. They preserve $\mathcal C$; this can be checked on its twelve basis words. In projective action $T^4=C$, and $C$ cycles $B_0,B_1,B_2$. All vector actions below are these coordinate permutations, not octonion-coordinate permutations.

\section{The smallest stable linear completion}
For $n=0,1,2,3$ retain $P(T^n\lambda)$. Their sums agree, so they give nine independent integer readouts, denoted $P_9$. Specifically, retain the first triple and the first two entries of each later triple.
\begin{proposition}\label{prop:ninefailure}
The kernel of $P_9$ is not $J$-stable. The linear span of its nine functionals and their $J$-translates has dimension 15. The smallest $\langle J,T\rangle$-stable real span containing them has dimension 24.
\end{proposition}
\begin{proof}
The lattice vector with numerator
\[
v=(8,-8,8,0,-8,0,-8,0,8,0,\ldots,0)
\]
has $P_9(v/\sqrt8)=0$ but
\[
P_9(Jv/\sqrt8)=(0,0,0,-2,-2,2,0,0,2).
\]
The code condition is trivial and the numerator sum is zero. Substitution verifies both readouts; rational row reduction gives ranks 9 and 15. For completeness of the stable span, act with $J,T$ on $B_0$, retaining every distinct resulting support. This bounded calculation produces exactly the 759 weight-eight words of $\mathcal C$. The next theorem proves that their indicator functionals span rank 24. This proves minimality, not just the existence of one 24-dimensional completion. The enumeration rule is complete: begin with $B_0$, apply both generators to every new support, and stop only when the list is closed. Both generators are invertible on the finite set.
\end{proof}

Let $\mathcal O$ be the set of all 759 octads. For an arbitrary real vector $x$, put
\[
 a_B=\frac14\sum_{i\in B}x_i,\quad
 S=\sum_{B\in\mathcal O}a_B,\quad
 S_i=\sum_{B\ni i}a_B.
\]
\begin{theorem}[Complete octad reconstruction]\label{thm:reconstruction}
The map $x\mapsto(a_B)_{B\in\mathcal O}$ is injective and has inverse on its image
\begin{equation}\label{eq:inverse}
\boxed{x_i=\frac{23S_i-7S}{1012}.}
\end{equation}
Its real image consists exactly of data satisfying
$a_B=\frac14\sum_{i\in B}x_i(a)$ for all $B$. This is an explicit set of linear equations, not an existence condition.
\end{theorem}
\begin{proof}
Let $A$ be the $759\times24$ indicator matrix. The Golay octads form the Steiner system $S(5,8,24)$ \cite[Theorem 1]{chapman}. Each coordinate lies in 253 octads and each pair in 77, by counting five-subsets. Hence
\[
A^{\mathsf T}A=176I+77\one\one^{\mathsf T}.
\]
In particular its rank is 24. Directly,
$S=253\sum x_i/4$ and $S_i=44x_i+77\sum x_j/4$; elimination yields \eqref{eq:inverse}. Applying the forward map tests precisely the displayed image equations. Both compositions are the identity on their stated domains.
\end{proof}

\section{Exactly one cyclic integral obstruction}
Let $L_{\mathcal O}$ be the integral lattice generated by $2\one_B/\sqrt8$, for all $B\in\mathcal O$.
\begin{lemma}\label{lem:indexfour}
\[
L_{\mathcal O}=\{x/\sqrt8\in\La:\sum_i x_i\equiv0\pmod{16}\},\qquad
[\La:L_{\mathcal O}]=4.
\]
Moreover, with $\upsilon=\one/(2\sqrt8)$,
\[
L_{\mathcal O}^*=\La+\Z\upsilon,\qquad
L_{\mathcal O}^*/\La\cong\Z/4\Z.
\]
\end{lemma}
\begin{proof}
Every generator satisfies the displayed congruence. On $\La$ the sum is divisible by four, and $x\mapsto\sum x_i/4\bmod4$ is onto: the vector $(-3,1,\ldots,1)/\sqrt8$ maps to 1. Thus the congruence sublattice has index four.

Here is a small, fully specified determinant certificate for the reverse inclusion. Sort the octad masks as increasing binary integers, numbered from zero. The 24 rows at positions
\[
\begin{gathered}
197,555,724,671,582,487,580,464,333,59,49,397,\\
55,427,249,304,329,72,195,630,672,16,180,402
\end{gathered}
\]
have determinant $2^{14}$. This is checked by exact fraction-free elimination. After scaling every row by $2/\sqrt8$, their covolume is $2^{14}/2^{12}=4$. They generate a sublattice of the congruence sublattice with the same covolume, so equality holds.

Since $\La$ is integral, $\La\subset L_{\mathcal O}^*$. Also $(\upsilon,2\one_B/\sqrt8)=1$, so $\upsilon\in L_{\mathcal O}^*$. The vector $4\upsilon$ lies in $\La$, but $2\upsilon$ does not, by \eqref{eq:model}. The quotient $L_{\mathcal O}^*/\La$ has order four by unimodularity and the index just proved. Hence $\upsilon$ generates it.
\end{proof}

\begin{theorem}[Integral gluing syndrome]\label{thm:syndrome}
Suppose $a\in\Z^{\mathcal O}$ satisfies the real image equations of Theorem~\ref{thm:reconstruction}. Then
\[
\delta(a)=\frac{11S-23S_0}{506}
\]
is an integer. There is a unique $\lambda=x/\sqrt8\in\La$ having these octad traces if and only if
\begin{equation}\label{eq:syndrome}
\boxed{\delta(a)\equiv0\pmod4.}
\end{equation}
When it exists, \eqref{eq:inverse} reconstructs all its coordinates. For example, the constant trace assignment $a_B=1$ is linearly compatible but has $\delta=5$, so it is not a Leech trace packet.
\end{theorem}
\begin{proof}
Compatibility gives a real vector pairing integrally with all generators of $L_{\mathcal O}$, hence an element of $L_{\mathcal O}^*=\La+\Z\upsilon$. Set $v_0=(-3,1,\ldots,1)/\sqrt8\in\La$. The homomorphism $z\mapsto4(z,v_0)\bmod4$ vanishes on $\La$ and sends $\upsilon$ to $5\bmod4$, so it identifies that order-four quotient. In raw coordinates it is
\[
4(x/\sqrt8,v_0)=(\sum_i x_i-4x_0)/2
=(11S-23S_0)/506.
\]
This proves integrality, the criterion, and uniqueness. Constant traces reconstruct $x_i=1/2$, that is $\upsilon$; their syndrome is five.
\end{proof}

\section{Scope and verification}
The output is a complete integral coordinate theorem and a minimal linear symmetry completion. It does not assert that the additional Erd\H{o}s--Straus cubic contains a positive integral trace. The data and maps are classical-code constructions; no new lattice, uniqueness theorem, or historical-priority claim is made. A content-level search of Golay design lattices, Smith invariants, octad trios and octonionic models did not establish a prior exact statement of this combined marked reconstruction; it also cannot establish absence from the literature.

The standalone \texttt{verify\_audit.py} reconstructs the code, the generator orbit, all entries of $A^{\mathsf T}A$, the determinant minor, the kernel witness, and exact inverse/syndrome tests. The finite facts entering the proof have fixed bounds ($4096$ codewords, $759$ octads, a $24\times24$ determinant). No floating point, unbounded search, formal-prover run or independent reviewer is used. The general assertions follow from the arguments above, not from extrapolating these tests.
